\scanpage{36}
\begin{center}ii\end{center}

some property of semi-simple algebraic groups in all characteristics which takes the place of the full reducibility of representations which is valid only in characteristic $0$. I conjecture, in fact, that if a semi-simple algebraic group $G$ is represented in a vector space $V$, and if $V_{\unsure{0|o}}$ is an invariant subspace of codimension $1$, then \uline{for} \uline{some} $\uline{\nu}$ the invariant subspace of codimension $1$
\[
V_{\unsure{0|o}}\cdot S^{p^\nu-1}(V)\subset S^{p^\nu}(V)
\]
is complemented by an invariant $1$-dimensional subspace.$^*$ Nonetheless, this is unknown, and one consequence is that this book is divided fairly sharply into two halves. Although both parts are closely analogous, they are logically independent. One half consists in Chapters 1, 2, 4, and 5 which deal essentially, only with characteristic zero, and yield a construction of the moduli scheme for curves over $\mathbb Q$. The other half consists in Chapters 3, 6, and 7, which deal with the ``arithmetic case'', i.e. over $\operatorname{Spec}(\mathbb Z)$, and yield a construction of the moduli schemes for curves over $\mathbb Z$. From a standpoint of content, however, Chapters 1, 2, 3, and 4 deal with orbit space problems, while Chapters 5, 6, and 7 deal with moduli.

Another unsolved problem deals with rationality over a perfect but non-algebraically closed field $k$, and is related to the Godement conjecture. In the notation of Chapter 2, the problem is whether, when

\noindent\rule{2cm}{.4pt}

$^*$Here $S^kV$ stands for the $k^{\mathrm{th}}$ symmetric power of $V$.

\mspage{35}
